\section{企業型年金の開始}

\subsection{目的}

\begin{InyouJobunbox}
  \begin{itemize}
    \item 法第1条
  \end{itemize}
\end{InyouJobunbox}

\begin{fillblank}[年金1 2025 問題1(2)アイ]
  \noindent \textbf{◯DC法}
  \begin{lawstructure}
    \begin{lawarticle}{第一条}{目的}
      この法律は、少子高齢化の進展、（　　A　　）の多様化等の社会経済情勢の変化にかんがみ、
      個人又は事業主が拠出した資金を個人が（　　B　　）において運用の指図を行い、
      高齢期においてその結果に基づいた給付を受けることができるようにするため、
      確定拠出年金について必要な事項を定め、国民の高齢期における所得の確保に係る自主的な努力を支援し、
      もって公的年金の給付と相まって国民の生活の安定と福祉の向上に寄与することを目的とする。
    \end{lawarticle}
  \end{lawstructure}
\end{fillblank}

\begin{answerbox}
 \begin{enumerate}[A:]
    \item 高齢期の生活
    \item 自己の責任
  \end{enumerate}
\end{answerbox}



\begin{fillblank}[年金1 \ H27 \ 問題1(2)]

  
  \noindent \textbf{◯確定給付企業年金法}
  \begin{lawstructure}
    \begin{lawarticle}{第一条}{目的}
      この法律は、少子高齢化の進展、（　　A　　）の変化等の社会経済情勢の変化にかんがみ、
      事業主が従業員と給付の内容を約し、高齢期において従業員がその内容に基づいた給付を受けること
      ができるようにするため、確定給付企業年金について必要な事項を定め、
      国民の高齢期における所得の確保に係る（　　B　　）を支援し、
      もって（　　C　　）と相まって国民の生活の安定と福祉の向上に寄与することを目的とする。
    \end{lawarticle}
  \end{lawstructure}

      
  \noindent \textbf{◯確定拠出年金法}
  \begin{lawstructure}
    \begin{lawarticle}{第一条}{目的}
      この法律は、少子高齢化の進展、（　　D　　）の多様化等の社会経済情勢の変化にかんがみ、
      個人又は事業主が拠出した資金を個人が（　　E　　）において運用の指図を行い、
      高齢期においてその結果に基づいた給付を受けることができるようにするため、
      確定拠出年金について必要な事項を定め、国民の高齢期における所得の確保に係る（　　B　　）を支援し、
      もって（　　C　　）と相まって国民の生活の安定と福祉の向上に寄与することを目的とする。
    \end{lawarticle}
  \end{lawstructure}

\end{fillblank}

\begin{answerbox}
 \begin{enumerate}[A:]
    \item 産業構造
    \item 自主的な努力
    \item 公的年金の給付
    \item 高齢期の生活
    \item 自己の責任
  \end{enumerate}
\end{answerbox}

\newpage

\subsection{各種用語の定義}

\begin{InyouJobunbox}
  \begin{itemize}
    \item 法第2条
    \item 令第1条
    \item 則第1条
  \end{itemize}
\end{InyouJobunbox}

\begin{truefalse}[年金1 2019 問題1(3)\textcircled{1}]
\textbf{下線部の内容の正誤を判定し、誤っている場合は正しい内容に改めなさい。}
\\

企業型年金における「記録関連業務」は次の業務を指す。
\begin{enumerate}
  \item 企業型年金加入者および企業型年金運用指図者ならびに個人型年金加入者および個人型
年金運用指図者（以下「加入者等」と総称する。）の氏名、住所、個人別管理資産額その
他の加入者等に関する事項の記録、保存および通知
\item 加入者等が行った運用の指図の取りまとめおよびその内容の資産管理機関（企業型年金
を実施する事業主が確定拠出年金法第八条第一項の規定により締結した契約の相手方をい
う。）または国民年金基金連合会への通知
\item \underline{運用の方法の選定および加入者等に対する提示}

\end{enumerate}
\end{truefalse}

\begin{answerbox}
  \textbf{誤}

  給付を受ける権利の裁定
\end{answerbox}

\newpage

\subsection{実施にかかる手続き（労使合意、承認）}

\begin{InyouJobunbox}
  \begin{itemize}
    \item 法第3条
    \item 令第3条
    \item 則第2条〜第3条の2
  \end{itemize}
\end{InyouJobunbox}

出題例なし

\newpage

\subsection{規約の承認基準}


\begin{InyouJobunbox}
  \begin{itemize}
    \item 法第4条
    \item 令第6条
    \item 則第4条、第4条の3
    \item 解釈第1-4〜第1-11
  \end{itemize}
\end{InyouJobunbox}

\begin{fillblank}[数理人試験 2021 問題1 設問5（改題）]
    \noindent \textbf{◯DC法}
  \begin{lawstructure}
    \begin{lawarticle}{第四条}{承認の基準等}
      厚生労働大臣は、前条第一項の承認の申請があった場合において、当該申請に係る規約が次に掲げる要件に適合すると認めるときは、同項の承認をするものとする。
      \lawitem{一〜二　（略）}
      \lawitem{三}
      事業主掛金について、（　　A　　）その他これに類する方法により算定した額によることが定められていること。

      （以下略）
    \end{lawarticle}
  \end{lawstructure}
\end{fillblank}

\begin{answerbox}
  \begin{enumerate}[A:]
    \item 定額または給与に一定の率を乗ずる方法
  \end{enumerate}
\end{answerbox}

\begin{explanationbox}
  簡易企業型年金は2026年4月1日付で廃止されている。
\end{explanationbox}


\begin{fillblank}[年金1 2025 問題1(2)ウエ]
    \noindent \textbf{◯DC則}
  \begin{lawstructure}
    \begin{lawarticle}{第四条}{企業型年金の給付の額の算定方法の基準}
      令第五条第一号の年金として支給されるものの算定方法は、次の各号に掲げる給付の区分に応じ、当該各号に定める基準に適合するものとする。

      \lawitem{一}
      年金たる老齢給付金

      \hangedlawitem{イ〜ロ　（略）}{3}

      \hangedlawitem{ハ}{3}
      給付の額（ホ及びチの規定により算定される額を除く。）は、請求日の属する月又はヘの申出をした日の属する月
      の前月の末日における個人別管理資産額の二分の一に相当する額を超えず、かつ、（　　A　　）を下回
      らないものであること（請求日において、個人別管理資産（当該企業型年金に係るものに限る。以下この条におい
      て同じ。）について、保険又は共済の契約であって（　　B　　）を支給することを約したものに基づく保険料又は共済
      掛金の払込みによって運用の指図を行っているものに係る給付の額を除く。ニにおいて同じ。）。

      （以下略）

    \end{lawarticle}
  \end{lawstructure}
\end{fillblank}

\begin{answerbox}
  \begin{enumerate}[A:]
    \item 二十分の一に相当する額
    \item 終身年金
  \end{enumerate}
\end{answerbox}




\newpage

\subsection{規約の変更}

\begin{InyouJobunbox}
  \begin{itemize}
    \item 法第5条、第6条
    \item 則第5条〜第7条の2
  \end{itemize}
\end{InyouJobunbox}

\begin{truefalse}[年金1 H29 問題1(4)\textcircled{5}]
  次のＡ～Ｄの４つの企業型年金規約の変更のうち、確定拠出年金法施行規則第５条第１項に
定める軽微な変更に該当するものを選択しなさい。
\begin{enumerate}[A:]
  \item 企業型年金を実施する事業主の名称に係る変更
  \item 運用の指図に関する事項の変更
  \item 年金の支給予定期間の種類の追加に係る変更
  \item 企業型年金の事業年度に関する事項の変更
\end{enumerate}
\end{truefalse}

\begin{answerbox}
  A,C,D
\end{answerbox}

\newpage

\subsection{運営管理業務の委託等}

\begin{InyouJobunbox}
  \begin{itemize}
    \item 法第7条、第8条
    \item 令第7条〜第9条
    \item 則第8条
  \end{itemize}
\end{InyouJobunbox}

\begin{truefalse}[自作 2026/10/17]
  
  運営管理業務の委託を行なっている事業主は、\underline{三年ごと}に、
  運営管理業務の実施に関する評価を行い、運営管理業務の委託について検討を加え、
  必要があるときは、確定拠出年金運営管理機関の変更その他の必要な措置を講ずるように努めなければならない。

\end{truefalse}

\begin{answerbox}
  少なくとも5年ごと
\end{answerbox}

\newpage